%% Ocean Engineering — CAS double-column template
%% "Visual Similarity Is Not Enough: Domain-Adapted Synthetic Data
%%  for Maritime Vessel Detection"

\documentclass[a4paper,fleqn]{cas-dc}

\usepackage[authoryear,longnamesfirst]{natbib}
\usepackage{booktabs}
\usepackage{multirow}
\usepackage{graphicx}
\usepackage{amsmath}
\usepackage{hyperref}
\usepackage{xcolor}
\usepackage{adjustbox}

% Removed \todo command — no remaining TODOs

\begin{document}
\let\WriteBookmarks\relax
\def\floatpagepagefraction{1}
\def\textpagefraction{.001}

\shorttitle{Visual Similarity Is Not Enough}
\shortauthors{D.L. Freire et al.}

\title[mode=title]{Visual Similarity Is Not Enough: Domain-Adapted Synthetic Data for Maritime Vessel Detection}

\author[1]{Daniela L. Freire}[orcid=0000-0000-0000-0000]
\cormark[1]
\ead{daniela.freire@usp.br}
\credit{Conceptualization, Methodology, Software, Investigation, Writing -- Original Draft, Visualization}

\author[2]{Eduardo H. Teixeira}[orcid=00000-0001-8449-7541]
\ead{eduardot@gea.inatel.br}
\credit{Data Curation, Resources, Writing -- Review \& Editing}

\author[3]{Leandro A. S. Moreira}[orcid=0000-0003-3470-6944]
\ead{jauvane@acm.org}
\credit{Supervision, Project Administration, Funding Acquisition, Writing -- Review \& Editing}

\affiliation[1]{organization={Institute of Mathematics and Computer Science (ICMC), University of S\~ao Paulo (USP)},
            city={S\~ao Carlos},
            state={SP},
            country={Brazil}}

\affiliation[2]{organization={Federal University of Itajubá (UNIFEI)},
            city={Itajubá},
            state={MG},
            country={Brazil}}

\affiliation[3]{organization={National Laboratory for Scientific Computing (LNCC)},
            city={Rio de Janeiro},
            state={RJ},
            country={Brazil}}

\cortext[1]{Corresponding author}

% ================================================================
% ABSTRACT
% ================================================================
\begin{abstract}
Operational maritime vessel detection systems require large annotated datasets for training, yet surveillance data from coastal and port-monitoring systems is often scarce and difficult to annotate. Public ship image datasets offer a potential solution, but their direct use as pre-training data may not transfer effectively to operational conditions. This paper investigates how to leverage the InaTechShips dataset ($\sim$28,000 public ship images) to improve YOLOv11m-based vessel detection on the CITRA-3D-Real dataset (2,081 operational maritime surveillance images). We first show that direct transfer learning from InaTechShips causes \emph{negative transfer}, reducing mAP50 by 4.15 percentage points compared to a COCO-pretrained baseline, due to a structural domain gap in scale (ships occupy $\sim$80\% vs $\sim$3\% of the image), object density (1 vs 3.37 objects per image), and capture conditions. We then propose \emph{in-place synthetic composition}: ship crops segmented with SAM from InaTechShips are resized and placed at the exact positions of real vessels in operational scenes. When used as joint training data alongside real images in a balanced regime, this approach achieves mAP50 of 0.8451 $\pm$ 0.0033, improving over the COCO baseline (0.8351 $\pm$ 0.0020) by 1.0 percentage point with non-overlapping confidence intervals. Our results indicate that both domain-adapted data generation \emph{and} an appropriate training regime are necessary to benefit from public ship imagery in operational maritime surveillance.
\end{abstract}

% ================================================================
% HIGHLIGHTS
% ================================================================
\begin{highlights}
\item Public ship pre-training causes negative transfer in maritime detection.
\item In-place synthetic composition bridges the public-operational domain gap.
\item Joint balanced training outperforms COCO baseline by 1.0 pp mAP50.
\item Training regime matters as much as data generation for transfer.
\item CLIP similarity fails to predict transfer for vessel detection.
\end{highlights}

% ================================================================
% KEYWORDS
% ================================================================
\begin{keywords}
Maritime vessel detection \sep Synthetic data augmentation \sep
Copy-paste augmentation \sep Domain adaptation \sep
Transfer learning \sep YOLO \sep Small object detection
\end{keywords}

\maketitle

% ================================================================
% 1. INTRODUCTION
% ================================================================
\section{Introduction}\label{sec:intro}

Maritime vessel detection is essential for coastal surveillance, search and rescue, port safety, maritime traffic monitoring, and autonomous surface vessel navigation. Deep learning approaches, particularly those based on the YOLO family of detectors \citep{redmon2016you, jocher2023ultralytics}, have achieved remarkable performance on standard benchmarks. However, deploying these models in operational maritime scenarios presents a fundamental challenge: operational datasets are typically small, restricted, and captured under conditions that differ substantially from publicly available maritime imagery.

The CITRA-3D-Real dataset, collected under an institutional maritime monitoring programme, exemplifies this challenge. With only 2,081 images containing 7,003 annotated bounding boxes, it provides limited training data. Moreover, its characteristics diverge significantly from public ship image datasets: vessels appear as small, distant objects (71.6\% are classified as ``small'' under COCO standards), multiple objects populate each scene (mean 3.37 per image), and images are captured under variable weather and lighting conditions from fixed surveillance cameras.

The InaTechShips dataset \citep{teixeira2025inatechships}, containing approximately 28,000 images sourced from the shipspotting.com community, offers an attractive source of additional training data. An intuitive hypothesis is that pre-training a detector on these public ship images---particularly those selected for visual similarity to the operational domain via CLIP embeddings \citep{radford2021learning}---would improve detection performance when fine-tuned on the operational dataset.

\textbf{This paper shows that this intuition does not hold in the studied operational setting, and proposes a solution.} Our investigation follows three stages:

\begin{enumerate}
    \item \textbf{Negative transfer discovery.} We show that pre-training on InaTechShips images---whether curated by CLIP similarity or randomly selected---degrades detection performance by 3.5--4.2\% mAP50 compared to using standard COCO pre-trained weights. An epoch ablation study reveals monotonic degradation: even 10 epochs of exposure to InaTechShips data reduces performance.
    
    \item \textbf{Domain gap diagnosis.} We identify three axes of structural incompatibility: scale (ships occupy $\sim$80\% of the image in InaTechShips vs $\sim$3\% in CITRA-3D-Real), object density (1 vs 3.37 per image), and capture context (professional photography vs operational surveillance).
    
    \item \textbf{In-place synthetic composition.} We propose a method that bridges the domain gap by extracting ship crops from InaTechShips using SAM \citep{kirillov2023segment}, resizing them to match the scale of real vessels, and placing them at the exact positions of annotated objects in operational scenes. This produces synthetic training images that preserve the spatial distribution, density, and visual context of the operational domain while introducing vessel appearance diversity from the public dataset.
\end{enumerate}

The main contributions of this work are:

\begin{itemize}
    \item Empirical evidence that CLIP-based visual similarity alone appears insufficient as a proxy for transfer learning compatibility in maritime object detection, with a controlled ablation isolating the effect of domain incompatibility from curation strategy.
    \item A domain adaptation method---in-place synthetic composition---that transforms public ship images into training data compatible with the operational domain by preserving scale, density, and context distributions.
    \item A systematic comparison of training regimes (sequential pre-training, frozen backbone, joint training) showing that the training strategy is as important as the data generation method: joint balanced training with synthetic and real data improves mAP50 by +1.0 percentage point over the COCO baseline, while sequential pre-training with the same synthetic data causes forgetting.
    \item A comprehensive experimental analysis including hyperparameter validation via HPO, epoch ablation, and comparison of curated vs.\ random selection, providing a reproducible framework for maritime detection research. All code, configurations, and analysis scripts are publicly available.\footnote{\url{https://github.com/DanielaLFreire/maritime-synthetic-detection}}
\end{itemize}

% ================================================================
% 2. RELATED WORK
% ================================================================
\section{Related Work}\label{sec:related}

\subsection{Maritime vessel detection}

Deep learning methods have transformed maritime vessel detection in recent years. The YOLO family, from its original formulation \citep{redmon2016you} to the latest YOLOv11 \citep{jocher2023ultralytics}, has become the dominant paradigm due to its balance of speed and accuracy. \citet{liu2025yolov8resattnet} proposed YOLOv8-ResAttNet, integrating residual attention mechanisms to improve detection in complex ocean backgrounds with wave reflections and cloud cover. \citet{li2024s3det} introduced S$^3$Det, specifically targeting small-scale ship detection by combining Normalized Wasserstein Distance loss with Feedback Cut\&Paste augmentation.

A persistent challenge is the detection of small and distant vessels. The SeaShips dataset \citep{shao2018seaships}, one of the largest precisely annotated maritime benchmarks, contains 31,455 images from coastal surveillance cameras. However, most public maritime datasets consist of close-range photographs where vessels are prominent, creating a mismatch with operational scenarios where ships appear as small, distant objects. This scale discrepancy is central to the domain gap investigated in this work.

Synthetic data generation has emerged as a promising approach to address data scarcity. \citet{ribeiro2022synthetic} demonstrated real-time ship segmentation using automatically annotated synthetic datasets, highlighting the potential of synthetic-to-real transfer in maritime environments.

\subsection{Copy-paste data augmentation}

The copy-paste augmentation strategy, first formalized by \citet{ghiasi2021simple}, creates synthetic training images by extracting object instances and compositing them onto different backgrounds. This approach has been particularly effective for instance segmentation and object detection tasks with limited data.

In the maritime domain, several recent works have applied copy-paste specifically to vessel detection. \citet{ruizponce2023poseidon} developed POSEIDON, a data augmentation tool for small object detection in maritime environments that uses metadata-guided sample generation. \citet{li2024s3det} proposed S$^3$Det with Feedback Cut\&Paste for small-scale ship detection, demonstrating improvements with YOLOv8. \citet{nemati2025maritime} applied controlled copy-paste with adaptive blending for maritime scenarios, validated with RT-DETR.

Our work differs from these approaches in a key aspect: rather than placing ships at random or heuristically-determined positions, we use the positions of real annotated vessels in operational images as anchors for synthetic composition, reducing the need for arbitrary placement decisions.

\subsection{Domain adaptation and transfer learning}

Transfer learning---initializing a detector with weights pre-trained on a large dataset such as COCO \citep{lin2014microsoft}---is standard practice in object detection. However, the assumption that domain-specific pre-training always improves downstream performance does not hold universally. \citet{zhang2023negative} survey the phenomenon of \emph{negative transfer}, where source domain knowledge degrades target task performance due to distributional mismatch between domains.

The problem is particularly acute in sequential fine-tuning, where \emph{catastrophic forgetting} \citep{mccloskey1989catastrophic, french1999catastrophic} causes the model to overwrite previously learned features. \citet{yang2024forgetting} provide a comprehensive taxonomy of forgetting phenomena in deep learning, distinguishing harmful forgetting (which degrades performance) from beneficial forgetting (which improves generalization). In the context of object detection, \citet{chen2018domain} proposed Domain Adaptive Faster R-CNN, using adversarial alignment at both image and instance levels to reduce domain shift. \citet{oza2023unsupervised} survey unsupervised domain adaptation methods for object detection, identifying feature alignment, pseudo-labeling, and image-level translation as the main strategies.

A complementary line of work questions the relationship between visual similarity and task-level transferability. \citet{zamir2018taskonomy} showed that task affinity does not follow intuitive expectations, and \citet{cui2018large} demonstrated that transfer learning effectiveness depends on the similarity of the underlying data distributions, not just visual appearance. Our work provides empirical evidence for this principle in the maritime detection domain: datasets that are visually similar (both containing ships) but distributionally incompatible (different scales, densities, contexts) produce negative rather than positive transfer.

% ================================================================
% 3. DATASETS AND EXPERIMENTAL SETUP
% ================================================================
\section{Datasets and Experimental Setup}\label{sec:setup}

Fig.~\ref{fig:composition_process} illustrates both the domain gap and the composition process. InaTechShips contains close-range photographs where vessels occupy $\sim$80\% of the frame (a), while CITRA-3D-Real captures distant vessels at operational scale (c). SAM segmentation isolates vessels with transparent backgrounds (b), which are then resized and composited at the exact positions of real vessels, producing the synthetic training image (d).

\begin{figure*}[!htbp]
  \centering
  \includegraphics[width=\textwidth]{figs/fig_composition_process.pdf}
  \caption{In-place synthetic composition pipeline. (a)~InaTechShips source image: close-range photograph from shipspotting.com, vessel occupying $\sim$80\% of the frame. (b)~SAM ViT-B segmentation using the bounding box as prompt, producing an RGBA crop with transparent background (checkerboard). (c)~CITRA-3D-Real operational scene with multiple distant vessels ($\sim$3\% of image area). (d)~Synthetic result: each real vessel is replaced by a randomly selected crop resized to the original bounding box dimensions and composited via alpha blending. The process generates 13 variations per scene, producing 27,053 synthetic images with 91,035 objects (3.37~obj/img). Labels are inherited from the operational annotations. Synthetic partitions are generated exclusively from the corresponding source partition to prevent data leakage.}
  \label{fig:composition_process}
\end{figure*}

\subsection{CITRA-3D-Real: operational maritime dataset}\label{sec:citra3d}

The CITRA-3D-Real dataset was collected under a maritime surveillance programme coordinated by CASNAV (Centre for Naval Systems Analysis), using fixed coastal surveillance cameras positioned at elevated vantage points overlooking a bay area. Images have a native resolution of 1920$\times$1080 pixels and capture a variety of maritime conditions including clear skies, overcast weather, haze, and varying sea states. The dataset spans multiple months, providing temporal diversity in lighting and weather conditions.

The dataset contains 2,081 images with 7,003 annotated bounding boxes across 9 original vessel classes: Tug (dominant), Cargo, Tanker, Passenger (8 instances only), Supply, Ferry, Military, Fishing, and Other. Due to incompatible taxonomies between datasets and severe class imbalance, all classes are collapsed into a single ``vessel'' class for detection. Annotations were produced by a single trained operator using LabelImg, with subsequent quality review. Institutional restrictions prevent disclosure of exact camera specifications and coordinates.

Table~\ref{tab:citra3d_profile} summarizes the scale profile of the dataset. The predominance of small objects (71.6\% under COCO criteria) and the high object density (median 2, maximum 40 per image) characterize a challenging detection scenario. Fig.~\ref{fig:scale_dist} contrasts the bounding box area distributions of CITRA-3D-Real and InaTechShips, revealing a $\sim$500$\times$ difference in median object area.

\begin{table}[!htbp]
\caption{Scale profile of CITRA-3D-Real bounding boxes (normalized coordinates). Statistics computed over 7,003 annotations in 2,081 images.}\label{tab:citra3d_profile}
\begin{adjustbox}{max width=\columnwidth}
\begin{tabular}{@{}lcccccc@{}}
\toprule
Metric & Mean & Median & P10 & P90 & Min & Max \\
\midrule
Width  & 0.047 & 0.031 & 0.011 & 0.102 & 0.004 & 0.551 \\
Height & 0.046 & 0.032 & 0.013 & 0.096 & 0.004 & 0.438 \\
Area   & 0.004 & 0.001 & 0.0002 & 0.008 & 0.00003 & 0.211 \\
Objects/img & 3.37 & 2.0 & 1.0 & 7.0 & 1 & 40 \\
\bottomrule
\end{tabular}
\end{adjustbox}
\end{table}

The dataset is split into training (1,348 images, 4,489 boxes), validation (332 images, 1,267 boxes), and test (401 images, 1,247 boxes) sets.

\begin{figure*}[!htbp]
  \centering
  \includegraphics[width=0.8\textwidth]{figs/fig3_scale_distribution.pdf}
  \caption{Distribution of normalized bounding box areas. (a)~CITRA-3D-Real: median area 0.10\% of image, 71.6\% classified as ``small'' under COCO criteria. (b)~InaTechShips: median area 54\% of image. The two distributions are nearly non-overlapping, reflecting a $\sim$500$\times$ difference in median object area.}
  \label{fig:scale_dist}
\end{figure*}

\subsection{InaTechShips: public ship image dataset}\label{sec:inatechships}

The InaTechShips dataset \citep{teixeira2025inatechships} is sourced from shipspotting.com, a community of maritime photography enthusiasts. Images are professional-quality photographs of individual vessels, typically captured from close range at ports and harbors. Ships occupy a large portion of the image ($\sim$60--80\%), with one object per image and controlled lighting conditions.

Two subsets are used in this work:

\begin{itemize}
    \item \textbf{dataset\_25k\_v2 (curated):} 27,796 images selected via CLIP ViT-B-32 cosine similarity $\geq$ 0.60 against CITRA-3D-Real embeddings. Split 60/20/20 with disjoint partitions.
    \item \textbf{random\_pool\_v2 (random control):} 27,964 images randomly sampled from the same source with stratified downsampling to match the temporal distribution of the curated subset.
\end{itemize}

Labels for both subsets were generated using PointRend \citep{kirillov2020pointrend} by the original InaTechShips authors, ensuring methodological symmetry.

\subsection{Detector and hyperparameters}\label{sec:detector}

All experiments use YOLOv11m \citep{jocher2023ultralytics} with single-class detection. Hyperparameters were validated through formal HPO using Optuna TPE \citep{akiba2019optuna} with 30 trials across 5 dimensions (learning rate, optimizer momentum, weight decay, warmup epochs, learning rate schedule). The best configuration found was within 1$\sigma$ of the baseline ($\Delta$mAP50 = $-$0.0023), leading to the decision to retain the baseline configuration: AdamW optimizer, $lr_0 = 0.001$, $lr_f = 0.01$, cosine learning rate schedule, momentum 0.937, weight decay 0.0005, 3 warmup epochs, image size 640, batch size 16.

\subsection{Evaluation protocol}\label{sec:protocol}

All experimental arms follow the same two-phase protocol: (1) pre-training on the source dataset for 100 epochs with patience 20, and (2) fine-tuning on CITRA-3D-Real for 300 epochs with patience 30. Evaluation is performed on the CITRA-3D-Real test set using mAP50 and mAP50-95 as primary metrics. Unless otherwise noted, experiments are repeated with three seeds (42, 123, 2024) and results are reported as mean $\pm$ standard deviation. Table~\ref{tab:protocol} summarises the experimental design.

\begin{table*}[!htbp]
\caption{Experimental protocol for each arm. All arms share the same hyperparameters and fine-tuning dataset; only the intermediate pre-training dataset varies.}\label{tab:protocol}
\begin{tabular}{@{}llclclc@{}}
\toprule
Arm & Init weights & Pre-train data & Epochs & Fine-tune data & Epochs & Seeds \\
\midrule
B1 & Random & --- & --- & CITRA-3D-Real & 300 & 3 \\
B2 & COCO & --- & --- & CITRA-3D-Real & 300 & 3 \\
A  & COCO & InaTechShips curated & 100 & CITRA-3D-Real & 300 & 3 \\
B  & COCO & InaTechShips random  & 100 & CITRA-3D-Real & 300 & 3 \\
A' & COCO & Synthetic (in-place) & 100 & CITRA-3D-Real & 300 & 3 \\
\bottomrule
\end{tabular}
\end{table*}

% ================================================================
% 4. METHOD: IN-PLACE SYNTHETIC COMPOSITION
% ================================================================
\section{Method: In-Place Synthetic Composition}\label{sec:method}

\subsection{Motivation}\label{sec:motivation}

As detailed in Section~\ref{sec:negative_transfer}, direct pre-training on InaTechShips images causes negative transfer due to a structural domain gap. Rather than discarding the public data, we propose adapting it to match the operational domain through synthetic image composition.

The key insight is that the domain gap manifests primarily in \emph{how} ships appear in the image (scale, position, context), not in \emph{what} ships look like (hull shape, superstructure). By extracting ship appearances from InaTechShips and placing them into the spatial context of CITRA-3D-Real, we can leverage the vessel diversity of the public dataset while respecting the operational characteristics of the target domain.

\subsection{Pipeline overview}\label{sec:pipeline}

The synthetic generation pipeline consists of four steps, illustrated visually in Fig.~\ref{fig:composition_process}:

\subsubsection{Step 1: Ship crop extraction with SAM}

For each image in the curated InaTechShips subset, we use the Segment Anything Model (SAM) \citep{kirillov2023segment} with ViT-B backbone to extract a precise segmentation mask of the vessel, using the annotated bounding box as a prompt. The resulting crop is saved as an RGBA image with transparent background, enabling seamless compositing. A quality filter retains crops with mask coverage between 25--95\% of the bounding box area, minimum dimension $\geq$ 50 pixels, and aspect ratio between 0.2--8.0, yielding 23,828 usable crops from 27,796 images (85.7\% yield).

\subsubsection{Step 2: In-place substitution}

For each CITRA-3D-Real image, we generate multiple variations by substituting each annotated vessel with a randomly selected crop from the filtered pool. The crop is resized to exactly match the dimensions of the original bounding box and composited using alpha blending. Crucially, the labels remain identical to the originals---the bounding box coordinates are inherited directly from the operational annotations.

This approach eliminates three sources of error present in conventional copy-paste methods:
\begin{itemize}
    \item \textbf{Placement errors:} ships are positioned exactly where real vessels were observed, never on land, buildings, or sky.
    \item \textbf{Scale errors:} ship sizes match the empirical distribution of the operational dataset by construction.
    \item \textbf{Density errors:} the number of objects per image matches the operational distribution exactly.
\end{itemize}

Note that at operational scales---where 71.6\% of vessels occupy less than 32$\times$32 pixels---the resized crop fully occludes the original vessel beneath it. For the 2.9\% of large objects, minor blending artifacts may occur at mask boundaries, but these did not measurably affect detection performance.

\subsubsection{Step 3: Dataset generation}

To prevent data leakage, the CITRA-3D-Real train/validation/test split was established before synthetic generation. Synthetic images for each partition were generated exclusively from the corresponding CITRA-3D-Real partition: training backgrounds produced only training synthetics, validation backgrounds produced only validation synthetics, and test backgrounds produced only test synthetics. No background, annotation, bounding box, or object position from the validation or test partitions was used during pre-training or model selection. Final evaluation was conducted solely on the held-out real CITRA-3D-Real test set.

We generate 13 variations per CITRA-3D-Real image, each with different randomly selected crops, producing a synthetic dataset of 27,053 images with 91,035 annotated objects (mean 3.37 per image, matching the operational distribution). The split mirrors the source: training (17,524 images), validation (4,316), and test (5,213).

% ================================================================
% 5. EXPERIMENTS AND RESULTS
% ================================================================
\section{Experiments and Results}\label{sec:results}

\subsection{Baselines}\label{sec:baselines}

Table~\ref{tab:baselines} shows the baseline results. COCO pre-training (B2) provides a substantial improvement over random initialization (B1): +4.3\% in mAP50 and +6.6\% in mAP50-95, with non-overlapping confidence intervals across the tested seeds.

\begin{table}[!htbp]
\caption{Baseline results on CITRA-3D-Real test set (3 seeds).}\label{tab:baselines}
\begin{adjustbox}{max width=\columnwidth}
\begin{tabular}{@{}llcc@{}}
\toprule
Arm & Pre-training & mAP50 & mAP50-95 \\
\midrule
B1 & Random init & 0.8008 $\pm$ 0.0061 & 0.4742 $\pm$ 0.0006 \\
\textbf{B2} & \textbf{COCO} & \textbf{0.8351 $\pm$ 0.0020} & \textbf{0.5055 $\pm$ 0.0022} \\
\bottomrule
\end{tabular}
\end{adjustbox}
\end{table}

\subsection{Negative transfer from direct pre-training}\label{sec:negative_transfer}

Table~\ref{tab:negative_transfer} presents the central finding: pre-training on InaTechShips images \emph{hurts} detection performance regardless of curation strategy.

\begin{table*}[!htbp]
\caption{Direct pre-training results. Both curated (CLIP) and random subsets cause negative transfer compared to B2.}\label{tab:negative_transfer}
\begin{tabular}{@{}llccc@{}}
\toprule
Arm & Pre-training data & mAP50 & mAP50-95 & $\Delta$ vs B2 \\
\midrule
B2 & COCO (baseline) & 0.8351 $\pm$ 0.0020 & 0.5055 $\pm$ 0.0022 & ref \\
A & Curated (3 seeds) & 0.7936 $\pm$ 0.0049 & 0.4692 $\pm$ 0.0017 & $-$4.15\% \\
B & Random (3 seeds) & 0.7945 $\pm$ 0.0046 & 0.4728 $\pm$ 0.0045 & $-$4.06\% \\
\bottomrule
\end{tabular}
\end{table*}

The curated and random subsets perform equivalently ($\Delta$ = 0.09 percentage points, within 1$\sigma$ of both arms' variance), confirming that the negative transfer is independent of the selection strategy and stems from structural domain incompatibility.

\subsection{Epoch ablation: pre-training duration analysis}\label{sec:ablation}

To distinguish between recoverable overfitting and fundamental incompatibility, we varied the number of pre-training epochs (Table~\ref{tab:ablation}, Fig.~\ref{fig:ablation}).

\begin{table}[!htbp]
\caption{Effect of pre-training duration on fine-tuning performance (seed 42). Degradation is monotonic.}\label{tab:ablation}
\begin{adjustbox}{max width=\columnwidth}
\begin{tabular}{@{}cccc@{}}
\toprule
Pre-training epochs & mAP50 & mAP50-95 & $\Delta$ mAP50 vs B2 \\
\midrule
0 (B2 baseline) & 0.8351 & 0.5055 & ref \\
10 & 0.8200 & 0.4999 & $-$1.51\% \\
20 & 0.8171 & 0.4960 & $-$1.80\% \\
50 & 0.8037 & 0.4731 & $-$3.14\% \\
100 (Arm A) & 0.8006 & 0.4680 & $-$3.45\% \\
\bottomrule
\end{tabular}
\end{adjustbox}
\end{table}

The degradation is strictly monotonic: each increment of pre-training worsens the final result. Even 10 epochs of exposure reduces mAP50 by 1.5\%. At 100 epochs, performance converges to the random initialization level (B1 = 0.8008), indicating that extended pre-training on InaTechShips completely erases the benefit of COCO weights.

\begin{figure}[!htbp]
  \centering
  \includegraphics[width=0.8\columnwidth]{figs/fig5_ablation_curve.pdf}
  \caption{Effect of pre-training duration on InaTechShips on downstream detection performance (CITRA-3D-Real test set, seed~42). Both mAP50 and mAP50-95 degrade monotonically with increasing pre-training epochs. Dashed lines indicate B2 (COCO baseline) and B1 (random initialization) reference levels.}
  \label{fig:ablation}
\end{figure}

\subsection{In-place synthetic composition results}\label{sec:synthetic_results}

The synthetic images generated by in-place composition were evaluated under multiple training regimes to identify the most effective strategy. Table~\ref{tab:main_results} and Fig.~\ref{fig:comparison} present the complete comparison.

\begin{table*}[!htbp]
\caption{Complete experimental results on CITRA-3D-Real test set. Joint balanced training with in-place synthetic data surpasses the COCO baseline. Sequential pre-training with 100 epochs causes forgetting even with domain-adapted data; shorter pre-training (20 epochs) or frozen backbone neutralise the gap.}\label{tab:main_results}
\begin{tabular}{@{}llcccccc@{}}
\toprule
Arm & Training strategy & mAP50 & mAP50-95 & P & R & F1 & Seeds \\
\midrule
\textbf{Joint balanced} & \textbf{COCO $\rightarrow$ real + synth (50/50)} & \textbf{0.8451 $\pm$ 0.0033} & \textbf{0.5206 $\pm$ 0.0017} & \textbf{0.857} & \textbf{0.805} & \textbf{0.830} & 3 \\
B2 (baseline) & COCO $\rightarrow$ CITRA-3D & 0.8351 $\pm$ 0.0020 & 0.5055 $\pm$ 0.0022 & 0.857 & 0.783 & 0.818 & 3 \\
Synth 20ep & COCO $\rightarrow$ synth (20ep) $\rightarrow$ CITRA-3D & 0.8344 $\pm$ 0.0033 & 0.5011 $\pm$ 0.0065 & 0.855 & 0.773 & 0.812 & 3 \\
Frozen backbone & COCO $\rightarrow$ synth (freeze) $\rightarrow$ CITRA-3D & 0.8342 $\pm$ 0.0039 & 0.5074 $\pm$ 0.0035 & 0.855 & 0.774 & 0.812 & 3 \\
Synth 100ep & COCO $\rightarrow$ synth (100ep) $\rightarrow$ CITRA-3D & 0.8221 $\pm$ 0.0085 & 0.4933 $\pm$ 0.0059 & 0.828 & 0.769 & 0.797 & 3 \\
B1 (baseline) & Random $\rightarrow$ CITRA-3D & 0.8008 $\pm$ 0.0061 & 0.4742 $\pm$ 0.0006 & 0.829 & 0.750 & 0.787 & 3 \\
B (random) & COCO $\rightarrow$ random pool $\rightarrow$ CITRA-3D & 0.7945 $\pm$ 0.0046 & 0.4728 $\pm$ 0.0045 & 0.858 & 0.742 & 0.796 & 3 \\
A (curated) & COCO $\rightarrow$ InaTechShips $\rightarrow$ CITRA-3D & 0.7936 $\pm$ 0.0049 & 0.4692 $\pm$ 0.0017 & 0.834 & 0.735 & 0.781 & 3 \\
\bottomrule
\end{tabular}
\end{table*}

Three key findings emerge from Table~\ref{tab:main_results}:

\textbf{(1) Joint balanced training is the best strategy.} Training on a balanced mixture of real and synthetic images (50\%/50\%, achieved by oversampling the real images 13$\times$) yields mAP50 = 0.8451 $\pm$ 0.0033, surpassing B2 by +1.0 percentage point with non-overlapping confidence intervals: [0.8418, 0.8484] vs.\ [0.8331, 0.8371]. The improvement in Recall is particularly notable (+2.2 pp), indicating that the model detects more vessels without sacrificing Precision.

\textbf{(2) Sequential pre-training causes forgetting even with adapted data.} With 100 epochs of sequential pre-training on synthetic data, performance drops to 0.8221 ($-$1.3\% vs B2), indicating that catastrophic forgetting occurs even when the pre-training data is domain-adapted. However, limiting pre-training to 20 epochs (0.8344) or freezing the backbone during pre-training (0.8342) neutralises the forgetting, achieving performance equivalent to B2.

\textbf{(3) The training regime matters as much as the data.} The same synthetic dataset produces results ranging from 0.8221 (sequential 100ep, below B2) to 0.8451 (joint balanced, above B2), a spread of 2.3 percentage points determined solely by how the data is used during training.

\begin{figure*}[!htbp]
  \centering
  \includegraphics[width=0.9\textwidth]{figs/fig6_comparison_bar.pdf}
  \caption{Comparison of all experimental arms on CITRA-3D-Real test set. (a)~mAP50: Joint balanced training achieves 0.845, surpassing the COCO baseline (B2~=~0.835) by +1.0 pp. Sequential pre-training arms (A, B) fall below the random initialisation baseline (B1). (b)~mAP50-95: similar pattern with +3.0\% relative improvement for joint balanced over B2. Error bars indicate $\pm$1 standard deviation across 3 seeds.}
  \label{fig:comparison}
\end{figure*}

% ================================================================
% 6. DISCUSSION
% ================================================================
\section{Discussion}\label{sec:discussion}

\subsection{Why direct transfer fails}

The negative transfer observed in Arms A and B is consistent with catastrophic forgetting or harmful intermediate feature adaptation \citep{mccloskey1989catastrophic, zhang2023negative}: the intermediate pre-training on InaTechShips likely overwrites generic low-level features learned from COCO (edges, textures, shapes) with domain-specific features optimised for close-range ship photography. These specialised features do not transfer to the operational domain, where ships appear as small, distant objects in complex maritime scenes. Paradoxically, the generic COCO features---learned from 80 diverse object categories with no dedicated vessel class---appear more useful than 28,000 ship-specific images.

\subsection{Why in-place composition works}

The synthetic composition resolves all three axes of the domain gap simultaneously:

\begin{itemize}
    \item \textbf{Scale:} ship crops are resized from $\sim$80\% to $\sim$3\% of the image, matching the operational distribution. At this scale, fine-grained appearance differences between professional and operational photography become imperceptible.
    \item \textbf{Density:} multiple ships per image (mean 3.37) match the operational distribution, training the model to detect objects in cluttered scenes.
    \item \textbf{Context:} backgrounds are real operational scenes, providing authentic ocean surfaces, weather conditions, and coastal features.
\end{itemize}

However, resolving the domain gap in the data is necessary but not sufficient: the training regime determines whether the adapted data improves or degrades performance.

\subsection{The role of training regime}

Our results reveal that the same synthetic dataset can produce outcomes ranging from $-$1.3\% (sequential 100-epoch pre-training) to +1.0\% (joint balanced training) relative to the COCO baseline. This 2.3 percentage point spread is determined entirely by how the synthetic data is presented to the model during training.

Sequential pre-training---even with domain-adapted data---creates a forgetting-recovery dynamic: the model first adapts to synthetic images, partially losing COCO features, then must recover during fine-tuning on limited real data. Two strategies mitigate this: short pre-training (20 epochs) limits the exposure and preserves most COCO features, while frozen backbone pre-training protects low-level features entirely. Both achieve performance equivalent to B2 but do not surpass it.

Joint balanced training eliminates the forgetting problem entirely by presenting real and synthetic images in every batch. The model simultaneously learns from authentic operational examples and diverse synthetic compositions, never ``forgetting'' either source. The 13$\times$ oversampling of real images ensures balanced representation (50\%/50\%) despite the asymmetry in dataset sizes. The improvement in Recall (+2.2 pp over B2, from 0.783 to 0.805) with maintained Precision (0.857) suggests that synthetic diversity helps the model detect a wider variety of vessel appearances.

\subsection{CLIP similarity as a proxy for transferability}

The equivalence of Arms A and B (curated $\approx$ random, $\Delta$ = 0.09 pp) suggests that global CLIP cosine similarity, when used alone, may fail to capture the structural properties that determine transfer learning success. The three axes of incompatibility we identified---scale distribution, object density, and capture conditions---are not encoded in CLIP's embedding space, which operates at the holistic image level rather than at the object-detection level.

\subsection{Limitations}

Several limitations should be acknowledged. First, our evaluation uses a single operational dataset (CITRA-3D-Real) from one geographic location with specific camera configurations; generalisation to other maritime surveillance systems remains to be validated. Second, all experiments use a single detector architecture (YOLOv11m); the extent to which our findings transfer to other detectors such as RT-DETR or Faster R-CNN is unknown, though the domain gap analysis is architecture-independent. Third, while the synthetic images are visually convincing at operational scales (where ships occupy $\sim$3\% of the image), close inspection reveals compositing artefacts at the crop boundaries---this is unlikely to affect detection training but may matter for tasks requiring pixel-level precision. Fourth, the method requires annotated bounding boxes in the target domain to serve as anchor positions for composition, limiting its applicability to fully unsupervised scenarios. Fifth, the crop resizing does not enforce aspect-ratio matching between source crop and target bounding box, which may introduce geometric distortion for objects with extreme proportions.

% ================================================================
% 7. CONCLUSION
% ================================================================
\section{Conclusion}\label{sec:conclusion}

This paper investigated how to effectively leverage public ship images for improving vessel detection in operational maritime surveillance scenarios. We showed that direct use of InaTechShips images as pre-training data causes negative transfer ($-$4.15\% mAP50), attributed to a structural domain gap in scale, density, and capture conditions. This gap affects curated and random subsets equally, indicating that CLIP-based visual similarity alone appears insufficient as a proxy for transfer learning compatibility.

We proposed in-place synthetic composition as a domain adaptation strategy: ship crops segmented with SAM are placed at the exact positions and scales of real vessels in operational scenes. This method requires no generative models and inherits target-domain annotations by construction. Critically, we found that the training regime determines whether the adapted data helps or hurts: sequential pre-training with 100 epochs still causes forgetting ($-$1.3\%), while joint balanced training with 50\% real and 50\% synthetic data improves mAP50 by +1.0 percentage point over the COCO baseline with non-overlapping confidence intervals, achieving the highest Recall (0.805) and F1 (0.830) across all tested configurations.

Although the method was evaluated on a single operational maritime dataset, the results suggest that structurally aligned synthetic composition, combined with an appropriate training regime, may be a promising strategy for domains in which public source data differ from target operational imagery in scale, density, and context.

Future work includes extending the approach to multi-class detection, where class-specific composition strategies may be needed to preserve inter-class balance. Testing with alternative detector architectures (RT-DETR, Faster R-CNN) would assess the generality of both the negative transfer finding and the composition solution. The optimal number of synthetic variations per image (currently 13) and the real-to-synthetic ratio in joint training could be further optimised. Finally, combining in-place composition with complementary augmentation strategies such as mosaic, mixup, or weather simulation may yield further improvements.

% ================================================================
% ACKNOWLEDGMENTS
% ================================================================
\section*{CRediT authorship contribution statement}

\textbf{Daniela L. Freire:} Conceptualisation, Methodology, Software, Investigation, Writing -- Original Draft, Visualisation.
\textbf{Eduardo H. Teixeira:} Data Curation, Resources, Writing -- Review \& Editing.
\textbf{Leandro A. S. Moreira:} Supervision, Project Administration, Funding Acquisition, Writing -- Review \& Editing.

\section*{Declaration of competing interest}

The authors declare that they have no known competing financial interests or personal relationships that could have appeared to influence the work reported in this paper.

\section*{Funding}

This work was supported by the Brazilian Navy under the CASNAV/DMarSup project (Termo 66/2025).

\section*{Acknowledgements}

We thank CASNAV for providing the CITRA-3D-Real dataset and Eduardo H. Teixeira for sharing the complete PointRend label set for InaTechShips. Computational resources were provided by Google Colab Pro+ (NVIDIA A100 GPU).

\section*{Declaration of generative AI and AI-assisted technologies in the writing process}

During the preparation of this work, the authors used Claude (Anthropic) for assistance with code development, data analysis pipeline design, and manuscript organisation. The authors reviewed and edited all AI-generated outputs and take full responsibility for the content of the publication.

\section*{Data and code availability}

All source code, experiment configurations, and analysis scripts are publicly available at \url{https://github.com/DanielaLFreire/maritime-synthetic-detection}. The repository includes 19 scripts organized in five stages (data preparation, dataset acquisition, analysis, synthetic generation, and training), YAML configuration files for each experimental arm, and JSON reports from all pipeline outputs. The CITRA-3D-Real dataset is restricted under the CASNAV/DMarSup project (Termo 66/2025) and available upon request to the corresponding author. The InaTechShips dataset is publicly accessible through the original authors' repository \citep{teixeira2025inatechships}.

% ================================================================
% REFERENCES
% ================================================================
\clearpage
\bibliographystyle{cas-model2-names}
\bibliography{references}

\end{document}
